% Table: Baseline Consistency Comparison
% Addresses P0-2: Shows improvement is robust across baseline choices

\begin{table}[t]
\centering
\begin{threeparttable}
\caption{\textbf{Geometric consistency comparison across baselines.}
Energy distance for PROP-DIRECT and two observation-agnostic baselines
under the sufficient observation regime. Relative reduction computed as
$(d_E(\text{baseline}) - d_E(\text{PROP})) / d_E(\text{baseline}) \times 100\%$.
Lower energy distance indicates better cross-view posterior consistency.
Both baselines yield comparable absolute energy distance, confirming
that relative improvements are robust to baseline choice.}
\label{tab:baseline_comparison}
\begin{tabular}{llcc}
\toprule
\textbf{Split} & \textbf{Model} & \textbf{Energy Distance} & \textbf{Relative Reduction} \\
\midrule
Test   & PROP-DIRECT  & $2.89 \times 10^{-4}$ & --- \\
       & B-STDPROB    & $4.57 \times 10^{-4}$ & 36.8\% \\
       & B-INDEP      & $4.66 \times 10^{-4}$ & 38.0\% \\
\midrule
Lockbox & PROP-DIRECT  & $2.86 \times 10^{-4}$ & --- \\
        & B-STDPROB    & $4.64 \times 10^{-4}$ & 38.3\% \\
        & B-INDEP      & $4.67 \times 10^{-4}$ & 38.6\% \\
\bottomrule
\end{tabular}
\begin{tablenotes}
\footnotesize
\item Bootstrap 90\% confidence intervals for the primary comparison
(PROP-DIRECT vs B-INDEP): Test [35.1\%, 40.7\%], Lockbox [34.8\%, 42.4\%].
B-STDPROB and B-INDEP differ by less than 1\% in absolute energy distance,
indicating that the choice of observation-agnostic baseline does not
materially affect conclusions.
\end{tablenotes}
\end{threeparttable}
\end{table}
